\chapter*{Tóm tắt đề tài}
\addcontentsline{toc}{chapter}{Tóm tắt đề tài}
SourceVerify là một công cụ web thử nghiệm giúp kiểm tra xem một ảnh có khả năng do AI tạo ra hay không. Thay vì dùng một mô hình đoán duy nhất, em cho nhiều phương pháp cùng phân tích ảnh rồi cộng điểm lại để ra kết quả cuối. Trong phạm vi Project I, em chỉ đi sâu vào 5 phương pháp tiêu biểu: \textit{Metadata Analysis}, \textit{Noise Residual}, \textit{DCT Block Artifacts}, \textit{Chromatic Aberration} và \textit{Spectral Nyquist Analysis}.

Báo cáo trình bày bối cảnh và lý do chọn đề tài, cơ sở lý thuyết của 5 phương pháp, cách thiết kế hệ thống và kết quả khi chạy thử trên một tập ảnh nhỏ. Với mỗi phương pháp em cố gắng nói rõ ý tưởng, cách tính điểm, ưu điểm và hạn chế. Kết quả cho thấy SourceVerify hợp với mục tiêu Project I: hiểu được bài toán, biết cách tiếp cận và làm được một sản phẩm minh họa có giải thích, dù chưa thể kết luận chắc chắn trong mọi trường hợp.

\textbf{Từ khóa:} ảnh AI, digital image forensics, metadata, noise residual, DCT, chromatic aberration, phân tích tần số.

\tableofcontents
\listoffigures
\listoftables
\clearpage
